\makeatletter
\def\input@path{{./}{../}{../../}{../../../}{../../../../}{preamble/}{../preamble/}{../../preamble/}{../../../preamble/}{../../../../preamble/}}
\makeatother
\input{preamble_exams.tex}
\input{graphs/G11_Histograms.tex}

\title{Grade 11 -- Term 1 \\ C5, C6 and C7 Exam Solutions -- Version 2}
\author{}
\date{}

\begin{document}
\maketitle
\vspace{-2cm}
\setcounter{tocdepth}{1}
\setcounter{secnumdepth}{2}
\phantomsection\label{toc}
\tableofcontents

\section*{Learning objective}\vspace{-5mm}
Use continuous probability models and sample data to describe typical values,
variability, and interval probabilities, and interpret differences between
sample estimates and theoretical values.

\section*{Criteria assessment}\vspace{-5mm}
This exam assesses C5, C6 and C7 separately, with 6 points per criterion
(18 points in total). Each task shows its point allocation. Partial credit
may be awarded when the correct procedure is clearly shown but arithmetic
errors are made. Answers without the required procedure will not receive credit.
Failure to write your name on the task sheet or submitting it after the deadline
will result in a one-point deduction from each criterion.
For this exam, a minimum of 4 out of 6 points is required to pass each criterion.
Use unrounded intermediate values; round final statistics to three decimal
places and percentage errors to two decimal places where necessary.

\section*{Evaluated criteria}
\begin{enumerate}[label=\textbf{C\arabic*},start=5]
\item Calculates and interprets the mean, variance and standard deviation of continuous distributions using standard formulas and functions.
\item Estimates interval probabilities using relative frequencies from generated samples.
 and compares them with theoretical probabilities.
\item Calculates sample means, variances, and standard deviations for different sample sizes and data representations.
\end{enumerate}


\sectionScored{Two custom-print studios}{C5-\sml{2}}{sec:c5-1}
{
Two studios print custom posters, with each order taking between 3 and 15 hours
to complete. Both studios handle a wide variety of orders, but their production
processes differ.
Studio I specializes in large, detailed posters that require extensive
customization. As a result, longer completion times are more common, with
orders becoming progressively more likely to require additional production
time.
Studio D uses a more efficient printing process designed for quick orders.
Shorter completion times are more common, and orders become progressively
less likely to require additional production time.
Let $I$ and $D$ represent the respective completion times in hours.
\QuestionsTasks[sec:c5-1]{
{5}/{sol:c5-1a}/{Identify each distribution from the clues and calculate its mean, variance, and standard deviation using the corresponding formulas. \score{2}},
}
\phantomsection\label{sol:c5-1a}
\subsubsection*{C5.1(a) -- Linear models \score{2}}
The intended models are a linearly increasing distribution for studio I
and a linearly decreasing distribution for studio D. The contextual trends
alone do not specify linear densities uniquely; the following results apply
under these intended linear models.
Thus $I\sim\operatorname{Lin}_{\mathrm{inc}}(3,15)$ and
$D\sim\operatorname{Lin}_{\mathrm{dec}}(3,15)$, with $a=3$, $b=15$.
\[
E(I)=\frac{a+2b}{3}=\frac{3+2(15)}{3}=11.000\text{ h},\qquad
E(D)=\frac{2a+b}{3}=\frac{2(3)+15}{3}=7.000\text{ h}.
\]
\[
Var(I)=Var(D)=\frac{(b-a)^2}{18}=\frac{12^2}{18}=8.000\text{ h}^2,
\qquad\sigma_I=\sigma_D=\sqrt{8}\approx2.828\text{ h}.
\]

}

\sectionScored{Bookings and daily takings}{C5-\sml{2}}{sec:c5-2}
{
A camera repair shop's daily revenue varies due to many small, independent
factors, including customer demand and repair costs. Most days, revenue
is close to \$3200, while higher and lower amounts become increasingly
uncommon at similar rates. The standard deviation is \$200.
Let $R$ be daily revenue in dollars.
The shop receives online bookings independently at an average rate of
10 per hour. Bookings arrive at a steady average rate, and the time since
the last booking does not affect when the next arrives.
Let $T$ be the time between bookings in minutes.
\QuestionsTasks[sec:c5-2]{
{5}/{sol:c5-2a}/{Identify the distribution for T from the clues and use its formulas to find the mean, variance, and standard deviation in minutes. \score{1}},
{5}/{sol:c5-2b}/{Identify the revenue distribution from the clues. Use its formulas to find the mean, variance, and standard deviation in dollars. Interpret the mean. \score{1}}
}
\phantomsection\label{sol:c5-2a}
\subsubsection*{C5.2(a) -- Exponential waiting time \score{1}}
Independent arrivals at a steady average rate and memoryless waiting identify
an exponential distribution. The rate is $\lambda=10/60=1/6$ per minute,
so $T\sim\operatorname{Exp}(1/6)$.
\[
E(T)=\frac{1}{\lambda}=6.000\text{ min},\qquad
Var(T)=\frac{1}{\lambda^2}=36.000\text{ min}^2,\qquad
\sigma_T=\frac{1}{\lambda}=6.000\text{ min}.
\]
\phantomsection\label{sol:c5-2b}
\subsubsection*{C5.2(b) -- Normal revenue and units \score{1}}
The symmetric pattern and many small, independent influences suggest a normal
model. Its mean is \$3200 and the stated standard deviation is
$\sigma_R=200$ dollars.
Thus $R\sim N(3200,40000)$ (mean, variance).
\[
E(R)=\mu_R=3200.000\text{ dollars},\qquad
Var(R)=\sigma_R^2=200^2=40000.000\text{ dollars}^2,
\]
\[
\sigma_R=\sqrt{Var(R)}=\sqrt{40000}=200.000\text{ dollars}.
\]
The long-run average daily revenue is \$3200; individual days may differ.
\secInicio[sec:c5-2]

}

\sectionScored{Two parcel-labeling stations}{C5-\sml{2}}{sec:c5-3}
{
A delivery company uses two parcel-labeling stations. At station A, an
automated machine labels parcels in 5 to 19 minutes. Due to differences
in parcel sizes and handling requirements, completion times vary
unpredictably, with no particular duration being more common than another.
Let $T_A$ be its labeling time in minutes.

At station B, staff label special parcels, also taking between 5 and
19 minutes. Most parcels are completed in around 10 minutes, while
shorter and longer processing times become progressively less common.
Very quick jobs are especially rare, and longer jobs occasionally
require additional handling. Let $T_B$ be its labeling time in minutes.

\QuestionsTasks[sec:c5-3]{
{5}/{sol:c5-3a}/{Identify the distribution for $T_A$ from the clues and calculate its mean, variance, and standard deviation using the corresponding formulas. \score{1}},
{5}/{sol:c5-3b}/{Identify the distribution for $T_B$ from the clues and calculate its mean, variance, and standard deviation using the corresponding formulas. \score{1}}
}
\phantomsection\label{sol:c5-3a}
\subsubsection*{C5.3(a) -- Uniform model \score{1}}
No particular duration being more common identifies the uniform model $T_A\sim U(5,19)$.
With $a=5$, $b=19$, apply $\mu=(a+b)/2$, $\sigma^2=(b-a)^2/12$,
and $\sigma=\sqrt{\sigma^2}$:
\[
E(T_A)=\frac{5+19}{2}=12.000\text{ min},\qquad
Var(T_A)=\frac{(19-5)^2}{12}=\frac{49}{3}\approx16.333\text{ min}^2,
\qquad\sigma_{T_A}=\sqrt{49/3}\approx4.041\text{ min}.
\]
\phantomsection\label{sol:c5-3b}
\subsubsection*{C5.3(b) -- Triangular model \score{1}}
The intended model is triangular, with bounds 5 and 19 and mode 10 minutes.
The contextual description alone does not uniquely specify a triangular density
or an exact mode; using this intended model gives
$T_B\sim\operatorname{Tri}(5,19,10)$
(minimum, maximum, mode). Apply $\mu=(a+b+c)/3$ and
$\sigma^2=(a^2+b^2+c^2-ab-ac-bc)/18$:
\[
E(T_B)=\frac{5+19+10}{3}=\frac{34}{3}\approx11.333\text{ min},
\]
\[
Var(T_B)=\frac{5^2+19^2+10^2-5(19)-5(10)-19(10)}{18}
=\frac{486-335}{18}=\frac{151}{18}\approx8.389\text{ min}^2,
\qquad\sigma_{T_B}=\sqrt{151/18}\approx2.896\text{ min}.
\]

}

\sectionScored{File upload times, grouped data}{C6-\sml{2}}{sec:c6-1}
{
A sample of file upload times in seconds is summarized below.
Classes are left-inclusive and right-exclusive.
\[
\begin{array}{c|rrrr}
\text{Class}&[2,8)&[8,14)&[14,20)&[20,26)\\\hline
f_i&5&9&7&3
\end{array}
\]
The theoretical probability is $P(8\le X<20)=0.75$.
\QuestionsTasks[sec:c6-1]{
{6}/{sol:c6-1a}/{Find the empirical probability of the stated interval from the given sample. \score{1}},
{6}/{sol:c6-1b}/{Calculate the absolute and percentage errors relative to the given theoretical probability. \score{1}}
}
\phantomsection\label{sol:c6-1a}
\subsubsection*{C6.1(a) -- Relative frequency \score{1}}
The event contains the two middle classes:
$n=5+9+7+3=24$ and $k=9+7=16$.
Because the event is a union of complete classes, no individual upload
times need to be estimated.
\[
\widehat P(8\le X<20)=\frac{16}{24}=\frac{2}{3}\approx0.667.
\]
\phantomsection\label{sol:c6-1b}
\subsubsection*{C6.1(b) -- Errors \score{1}}
Using the exact relative frequency,
\[
E=\left|\frac{2}{3}-0.75\right|=\frac{1}{12}\approx0.083,\qquad
\text{Percentage error}=\frac{\frac{1}{12}}{0.75}\times100\%\approx11.11\%.
\]
\secInicio[sec:c6-1]

}

\sectionScored{Office electricity use, raw data}{C6-\sml{1}}{sec:c6-2}
{
A sample of daily office electricity use over 20 working days (tens of kilowatt-hours) is
\[
\begin{array}{cccccccccc}
1.3&1.7&2.0&2.2&2.6&2.9&3.3&3.5&3.8&4.2\\
1.6&1.9&2.3&2.7&3.1&3.4&3.6&3.9&4.4&4.6
\end{array}
\]
Let $X$ be daily electricity use in tens of kilowatt-hours. The interval
20 to 35 kWh, including both limits, is $2\le X\le3.5$ in these units.
The theoretical probability is $P(2\le X\le3.5)=0.60$.
\QuestionsTasks[sec:c6-2]{
{6}/{sol:c6-2a}/{Find the empirical probability of the stated interval from the given sample. \score{0.5}},
{6}/{sol:c6-2b}/{Calculate the absolute and percentage errors relative to the given theoretical probability. \score{0.5}}
}
\phantomsection\label{sol:c6-2a}
\subsubsection*{C6.2(a) -- Relative frequency \score{0.5}}
There are $n=20$ observations. The included observations are
$2.0,2.2,2.6,2.9,3.3,3.5,2.3,2.7,3.1,3.4$, so $k=10$.
\[
\widehat P(2\le X\le3.5)=\frac{10}{20}=\frac{1}{2}\approx0.500.
\]
\phantomsection\label{sol:c6-2b}
\subsubsection*{C6.2(b) -- Errors \score{0.5}}
Using the exact relative frequency,
\[
E=\left|\frac{1}{2}-0.60\right|=\frac{1}{10}\approx0.100,\qquad
\text{Percentage error}=\frac{\frac{1}{10}}{0.60}\times100\%\approx16.67\%.
\]
\secInicio[sec:c6-2]

}

\sectionScored{Museum entrance queues, histogram data}{C6-\sml{2}}{sec:c6-3}
{
The histogram is the only record of a sample of museum entrance queue times.
Bar heights are frequencies; times are in minutes, with left-inclusive,
right-exclusive classes.
\begin{center}
\HistogramGraph{(1,5),(5,9),(9,13),(13,17)}{8,15,11,6}
{0.60\textwidth}{4.5cm}{blue!25}{4}{5}
\end{center}
The theoretical probability is $P(5\le X<13)=0.70$.
\QuestionsTasks[sec:c6-3]{
{6}/{sol:c6-3a}/{Find the empirical probability of the stated interval from the given sample. \score{1}},
{6}/{sol:c6-3b}/{Calculate the absolute and percentage errors relative to the given theoretical probability. \score{1}}
}
\phantomsection\label{sol:c6-3a}
\subsubsection*{C6.3(a) -- Relative frequency \score{1}}
The classes $[1,5),[5,9),[9,13),[13,17)$ have frequencies
$8,15,11,6$. Thus $n=8+15+11+6=40$ and $k=15+11=26$.
\[
\widehat P(5\le X<13)=\frac{26}{40}=\frac{13}{20}\approx0.650.
\]
\phantomsection\label{sol:c6-3b}
\subsubsection*{C6.3(b) -- Errors \score{1}}
Using the exact relative frequency,
\[
E=\left|\frac{13}{20}-0.70\right|=\frac{1}{20}\approx0.050,\qquad
\text{Percentage error}=\frac{\frac{1}{20}}{0.70}\times100\%\approx7.14\%.
\]
\secInicio[sec:c6-3]

}

\sectionScored{Equipment rental times, frequency data}{C6-\sml{1}}{sec:c6-4}
{
Recorded equipment rental times are summarized below.
\[
\begin{array}{c|rrrrrr}
x_i\text{ (min)}&4&8&12&16&20&24\\\hline
f_i&8&16&28&22&14&12
\end{array}
\]
The theoretical probability is $P(8\le X<20)=0.80$.
\QuestionsTasks[sec:c6-4]{
{6}/{sol:c6-4a}/{Find the empirical probability of the stated interval from the given sample. \score{0.5}},
{6}/{sol:c6-4b}/{Calculate the absolute and percentage errors relative to the given theoretical probability. \score{0.5}}
}
\phantomsection\label{sol:c6-4a}
\subsubsection*{C6.4(a) -- Relative frequency \score{0.5}}
The total is $n=8+16+28+22+14+12=100$.
The value 20 is excluded, so $k=16+28+22=66$.
\[
\widehat P(8\le X<20)=\frac{66}{100}=\frac{33}{50}\approx0.660.
\]
\phantomsection\label{sol:c6-4b}
\subsubsection*{C6.4(b) -- Errors \score{0.5}}
Using the exact relative frequency,
\[
E=\left|\frac{33}{50}-0.80\right|=\frac{7}{50}\approx0.140,\qquad
\text{Percentage error}=\frac{\frac{7}{50}}{0.80}\times100\%\approx17.50\%.
\]
\secInicio[sec:c6-4]

}

Use the sample-variance denominator $n-1$ throughout these four problems.
For each statistic, calculate absolute error $E=|\text{sample value}-\text{theoretical value}|$
and percentage error $100E/|\text{theoretical value}|$. Use unrounded values in all error calculations.
For grouped data and histograms, use class midpoints to estimate the statistics.

\sectionScored{Luggage masses, histogram data}{C7-\sml{2}}{sec:c7-1}
{
A histogram gives a sample of luggage masses in kilograms. Bar heights are
frequencies; classes are left-inclusive and right-exclusive.
The model has $\mu=9$ kg and $\sigma^2=7$ kg$^2$.
\begin{center}
\HistogramGraph{(4,7),(7,10),(10,13)}{3,4,2}
{0.55\textwidth}{4.5cm}{orange!30}{3}{2}
\end{center}
\QuestionsTasks[sec:c7-1]{
{7}/{sol:c7-1a}/{Using the grouped data, estimate the sample mean, sample variance, and sample standard deviation. \score{1}},
{7}/{sol:c7-1b}/{Calculate the absolute and percentage errors of the mean, variance, and standard deviation relative to the theoretical values. \score{1}}
}
\phantomsection\label{sol:c7-1a}
\subsubsection*{C7.1(a) -- Histogram statistics \score{1}}
The class midpoints are $5.5,8.5,11.5$ kg. Use these as the representative observations.
\[
n=3+4+2=9,\qquad
\bar{x}\approx\frac{3(\frac{11}{2})+4(\frac{17}{2})+2(\frac{23}{2})}{9}=\frac{49}{6}\approx8.167.
\]
The weighted squared-deviation sum, using the unrounded mean, is
\[
3\left(\frac{11}{2}-\frac{49}{6}\right)^2+4\left(\frac{17}{2}-\frac{49}{6}\right)^2+2\left(\frac{23}{2}-\frac{49}{6}\right)^2=44.
\]
\[
s^2\approx\frac{44}{9-1}=\frac{11}{2}\approx5.500,\qquad
s\approx\sqrt{\frac{11}{2}}\approx2.345.
\]
The mean and standard deviation are in kg; the variance is in squared kg.

\phantomsection\label{sol:c7-1b}
\subsubsection*{C7.1(b) -- Absolute and percentage errors \score{1}}
The theoretical standard deviation is $\sigma=\sqrt{7}$.
Apply $E=|\text{sample value}-\text{theoretical value}|$ and
$100E/|\text{theoretical value}|$ to the unrounded results:
\[
\begin{array}{l|rr}
\text{Statistic}&\text{Absolute error}&\text{Percentage error}\\\hline
Mean&0.833&9.26\%\\
Variance&1.500&21.43\%\\
Standard deviation&0.301&11.36\%\\
\end{array}
\]
Absolute errors have the units of the corresponding statistic; percentage errors are unitless.
These errors use midpoint estimates because individual observations are unknown.
\secInicio[sec:c7-1]

}

\sectionScored{Chamber temperatures, raw data}{C7-\sml{1}}{sec:c7-2}
{
Five chamber temperature readings in degrees Celsius are $14,17,18,19,22$.
The theoretical model has $\mu=19$ degrees Celsius and $\sigma^2=6$ square degrees Celsius.
\QuestionsTasks[sec:c7-2]{
{7}/{sol:c7-2a}/{Calculate the sample mean, sample variance, and sample standard deviation from the given data. \score{0.5}},
{7}/{sol:c7-2b}/{Calculate the absolute and percentage errors of the mean, variance, and standard deviation relative to the theoretical values. \score{0.5}}
}
\phantomsection\label{sol:c7-2a}
\subsubsection*{C7.2(a) -- Raw-data statistics \score{0.5}}
\[
n=1+1+1+1+1=5,\qquad
\bar{x}=\frac{1(14)+1(17)+1(18)+1(19)+1(22)}{5}=18\approx18.000.
\]
\[
\begin{array}{c|rrrrr}
x_i&14&17&18&19&22\\\hline
x_i-\bar{x}&-4&-1&0&1&4\\
(x_i-\bar{x})^2&16&1&0&1&16
\end{array}
\]
The squared deviations sum to $16+1+0+1+16=34$.
\[
s^2=\frac{34}{5-1}=\frac{17}{2}\approx8.500,\qquad
s=\sqrt{\frac{17}{2}}\approx2.915.
\]
The mean and standard deviation are in degrees Celsius; the variance is in squared degrees Celsius.

\phantomsection\label{sol:c7-2b}
\subsubsection*{C7.2(b) -- Absolute and percentage errors \score{0.5}}
The theoretical standard deviation is $\sigma=\sqrt{6}$.
Apply $E=|\text{sample value}-\text{theoretical value}|$ and
$100E/|\text{theoretical value}|$ to the unrounded results:
\[
\begin{array}{l|rr}
\text{Statistic}&\text{Absolute error}&\text{Percentage error}\\\hline
Mean&1.000&5.26\%\\
Variance&2.500&41.67\%\\
Standard deviation&0.466&19.02\%\\
\end{array}
\]
Absolute errors have the units of the corresponding statistic; percentage errors are unitless.
\secInicio[sec:c7-2]

}

\sectionScored{Parcels dispatched, frequency data}{C7-\sml{1}}{sec:c7-3}
{
The numbers of parcels dispatched in nine short collection periods are summarized below.
The theoretical model has $\mu=7$ parcels and $\sigma^2=9$ parcels$^2$.
\[
\begin{array}{c|rrr}x_i&3&7&11\\\hline f_i&2&4&3\end{array}
\]
\QuestionsTasks[sec:c7-3]{
{7}/{sol:c7-3a}/{Calculate the sample mean, sample variance, and sample standard deviation from the given data. \score{0.5}},
{7}/{sol:c7-3b}/{Calculate the absolute and percentage errors of the mean, variance, and standard deviation relative to the theoretical values. \score{0.5}}
}
\phantomsection\label{sol:c7-3a}
\subsubsection*{C7.3(a) -- Weighted statistics \score{0.5}}
\[
n=2+4+3=9,\qquad
\bar{x}=\frac{2(3)+4(7)+3(11)}{9}=\frac{67}{9}\approx7.444.
\]
The weighted squared-deviation sum, using the unrounded mean, is
\[
2\left(3-\frac{67}{9}\right)^2+4\left(7-\frac{67}{9}\right)^2+3\left(11-\frac{67}{9}\right)^2=\frac{704}{9}.
\]
\[
s^2=\frac{\frac{704}{9}}{9-1}=\frac{88}{9}\approx9.778,\qquad
s=\sqrt{\frac{88}{9}}\approx3.127.
\]
The mean and standard deviation are in parcels; the variance is in squared parcels.

\phantomsection\label{sol:c7-3b}
\subsubsection*{C7.3(b) -- Absolute and percentage errors \score{0.5}}
The theoretical standard deviation is $\sigma=\sqrt{9}$.
Apply $E=|\text{sample value}-\text{theoretical value}|$ and
$100E/|\text{theoretical value}|$ to the unrounded results:
\[
\begin{array}{l|rr}
\text{Statistic}&\text{Absolute error}&\text{Percentage error}\\\hline
Mean&0.444&6.35\%\\
Variance&0.778&8.64\%\\
Standard deviation&0.127&4.23\%\\
\end{array}
\]
Absolute errors have the units of the corresponding statistic; percentage errors are unitless.
\secInicio[sec:c7-3]

}

\sectionScored{Walking times, grouped data}{C7-\sml{2}}{sec:c7-4}
{
Walking times in minutes are grouped using left-inclusive, right-exclusive classes.
The theoretical mean is 23 minutes and variance is 30 minutes$^2$.
\[
\begin{array}{c|rrr}
\text{Time (min)}&[12,18)&[18,24)&[24,30)\\\hline
f_i&2&3&4
\end{array}
\]
\QuestionsTasks[sec:c7-4]{
{7}/{sol:c7-4a}/{Using the grouped data, estimate the sample mean, sample variance, and sample standard deviation. \score{1}},
{7}/{sol:c7-4b}/{Calculate the absolute and percentage errors of the mean, variance, and standard deviation relative to the theoretical values. \score{1}}
}
\phantomsection\label{sol:c7-4a}
\subsubsection*{C7.4(a) -- Midpoint estimates \score{1}}
The class midpoints are $15,21,27$ min. Use these as the representative observations.
\[
n=2+3+4=9,\qquad
\bar{x}\approx\frac{2(15)+3(21)+4(27)}{9}=\frac{67}{3}\approx22.333.
\]
The weighted squared-deviation sum, using the unrounded mean, is
\[
2\left(15-\frac{67}{3}\right)^2+3\left(21-\frac{67}{3}\right)^2+4\left(27-\frac{67}{3}\right)^2=200.
\]
\[
s^2\approx\frac{200}{9-1}=25\approx25.000,\qquad
s\approx\sqrt{25}\approx5.000.
\]
The mean and standard deviation are in min; the variance is in squared min.

\phantomsection\label{sol:c7-4b}
\subsubsection*{C7.4(b) -- Absolute and percentage errors \score{1}}
The theoretical standard deviation is $\sigma=\sqrt{30}$.
Apply $E=|\text{sample value}-\text{theoretical value}|$ and
$100E/|\text{theoretical value}|$ to the unrounded results:
\[
\begin{array}{l|rr}
\text{Statistic}&\text{Absolute error}&\text{Percentage error}\\\hline
Mean&0.667&2.90\%\\
Variance&5.000&16.67\%\\
Standard deviation&0.477&8.71\%\\
\end{array}
\]
Absolute errors have the units of the corresponding statistic; percentage errors are unitless.
These errors use midpoint estimates because individual observations are unknown.
\secInicio[sec:c7-4]


\newpage
\section{Summarized Solutions}
\begingroup
\small
\linespread{1.08}\selectfont
\setlength{\parskip}{5pt}
\setlength{\tabcolsep}{5pt}
\renewcommand{\arraystretch}{1.2}
\textit{Notation and rounding.} $\mu=E(X)$, $\sigma^2=\operatorname{Var}(X)$,
$\sigma$ = theoretical SD; $\bar{x},s^2,s$ = sample mean, variance, SD.
Decimal columns show three places (rounded where necessary); percentage errors
show two. Model parameters, counts and fractions are exact. All errors use
unrounded values.

\begin{minipage}{\linewidth}
\subsection*{C5}
\textbf{\ref{sec:c5-1}. Two custom-print studios}\\
\textbf{\ref{sec:c5-2}. Bookings and daily takings}\\
\textbf{\ref{sec:c5-3}. Two parcel-labeling stations}

\textbf{Distribution models and theoretical statistics}

\par\smallskip
\begin{tabularx}{\linewidth}{@{}l >{\raggedright\arraybackslash}X rrr l@{}}
\toprule
Question & Variable/model & $E(X)$ & $\operatorname{Var}(X)$ & $\sigma$ & Units \\
\midrule
\ref{sec:c5-1}, I & $I\sim\operatorname{Lin}_{\mathrm{inc}}(3,15)$ & 11.000 & 8.000 & 2.828 & h; h$^2$ \\
\ref{sec:c5-1}, D & $D\sim\operatorname{Lin}_{\mathrm{dec}}(3,15)$ & 7.000 & 8.000 & 2.828 & h; h$^2$ \\
\ref{sec:c5-2}, $T$ & $T\sim\operatorname{Exp}(1/6)$ & 6.000 & 36.000 & 6.000 & min; min$^2$ \\
\ref{sec:c5-2}, $R$ & $R\sim N(3200,40000)$ & 3200.000 & 40000.000 & 200.000 & dollars; dollars$^2$ \\
\ref{sec:c5-3}, A & $T_A\sim U(5,19)$ & 12.000 & 16.333 & 4.041 & min; min$^2$ \\
\ref{sec:c5-3}, B & $T_B\sim\operatorname{Tri}(5,19,10)$ & 11.333 & 8.389 & 2.896 & min; min$^2$ \\
\bottomrule
\end{tabularx}
\par\smallskip

\textbf{Mean and variance: substitute before calculating}

\par\smallskip
\begin{tabular*}{\linewidth}{@{\extracolsep{\fill}}l c c@{}}
\toprule
Question & $E(X)$: substituted expression & $\operatorname{Var}(X)$: substituted expression \\
\midrule
\ref{sec:c5-1}, I & $\dfrac{3+2(15)}{3}$ & $\dfrac{(15-3)^2}{18}$ \\
\ref{sec:c5-1}, D & $\dfrac{2(3)+15}{3}$ & $\dfrac{(15-3)^2}{18}$ \\
\ref{sec:c5-2}, $T$ & $\dfrac{1}{1/6}$ & $\dfrac{1}{(1/6)^2}$ \\
\ref{sec:c5-2}, $R$ & $\mu_R=3200$ & $200^2$ \\
\ref{sec:c5-3}, A & $\dfrac{5+19}{2}$ & $\dfrac{(19-5)^2}{12}$ \\
\ref{sec:c5-3}, B & $\dfrac{5+19+10}{3}$ & $\dfrac{5^2+19^2+10^2-5(19)-5(10)-19(10)}{18}$ \\
\bottomrule
\end{tabular*}
\par\smallskip

Units: first for $E(X),\sigma$; second for $\operatorname{Var}(X)$.
Conventions: $N(\mu,\sigma^2)$; $\operatorname{Tri}(a,b,c)$ = (minimum, maximum, mode);
exponential rate $\lambda=1/6\text{ min}^{-1}$.

\textit{Model notes.} Question \ref{sec:c5-1}: intended increasing/decreasing linear densities on $[3,15]$;
the contextual trends do not uniquely determine linear models.
Question \ref{sec:c5-3}: intended triangular density with mode 10 min; the description alone
specifies neither a unique triangular density nor an exact mode.
Question \ref{sec:c5-2}: \$3200 is the long-run average daily revenue; individual days may differ.
\end{minipage}

\begin{minipage}{\linewidth}
\subsection*{C6}
\textbf{\ref{sec:c6-1}. File upload times, grouped data}\\
\textbf{\ref{sec:c6-2}. Office electricity use, raw data}\\
\textbf{\ref{sec:c6-3}. Museum entrance queues, histogram data}\\
\textbf{\ref{sec:c6-4}. Equipment rental times, frequency data}

\textbf{Empirical probabilities and errors}

$P_0$ = theoretical probability, $\quad E=|\widehat P-P_0|,\quad
E\%=100E/|P_0|$.
For each question: $n,k,\widehat P$ and errors $E,E\%$.

\par\smallskip
\begin{tabularx}{\linewidth}{@{}l >{\raggedright\arraybackslash}X rr r r r r@{}}
\toprule
Question & Event & $n$ & $k$ & $\widehat P=k/n$ & $P_0$ & $E$ & $E\%$ \\
\midrule
\ref{sec:c6-1} & $8\le X<20$ s & 24 & 16 & $\frac{16}{24}\approx0.667$ & 0.750 & 0.083 & 11.11\% \\
\ref{sec:c6-2} & $2\le X\le3.5$ & 20 & 10 & $\frac{10}{20}=0.500$ & 0.600 & 0.100 & 16.67\% \\
\ref{sec:c6-3} & $5\le X<13$ min & 40 & 26 & $\frac{26}{40}=0.650$ & 0.700 & 0.050 & 7.14\% \\
\ref{sec:c6-4} & $8\le X<20$ min & 100 & 66 & $\frac{66}{100}=0.660$ & 0.800 & 0.140 & 17.50\% \\
\bottomrule
\end{tabularx}
\par\smallskip

Question \ref{sec:c6-2}: $X$ is in tens of kWh; $2\le X\le3.5$ means $20\le\text{use}\le35$ kWh.
Questions \ref{sec:c6-1} and \ref{sec:c6-3}: complete classes give exact counts; no midpoint estimates.
Probabilities and their errors are unitless.
\end{minipage}

\newpage
\begin{minipage}{\linewidth}
\subsection*{C7}
\textbf{\ref{sec:c7-1}. Luggage masses, histogram data}\\
\textbf{\ref{sec:c7-2}. Chamber temperatures, raw data}\\
\textbf{\ref{sec:c7-3}. Parcels dispatched, frequency data}\\
\textbf{\ref{sec:c7-4}. Walking times, grouped data}

\textbf{Sample statistics and comparison with theory}

$s^2=\dfrac{\sum_i f_i(x_i-\bar{x})^2}{n-1}$, with $f_i=1$ for raw data;
$\quad E=|\text{sample value}-\text{theoretical value}|$, $\quad
E\%=100E/|\text{theoretical value}|$.

\textbf{Sample statistics}

\par\smallskip
\begin{tabular*}{\linewidth}{@{\extracolsep{\fill}}l r rrr l@{}}
\toprule
Question & $n$ & $\bar{x}$ & $s^2$ & $s$ & Units \\
\midrule
\ref{sec:c7-1} & 9 & 8.167 & 5.500 & 2.345 & kg; kg$^2$ \\
\ref{sec:c7-2} & 5 & 18.000 & 8.500 & 2.915 & $^\circ$C; $(^\circ\mathrm{C})^2$ \\
\ref{sec:c7-3} & 9 & 7.444 & 9.778 & 3.127 & parcels; parcels$^2$ \\
\ref{sec:c7-4} & 9 & 22.333 & 25.000 & 5.000 & min; min$^2$ \\
\bottomrule
\end{tabular*}
\par\smallskip

Units: first for $\bar{x},s$; second for $s^2$.
\textit{Midpoint estimates:} Question \ref{sec:c7-1} uses $(5.5,8.5,11.5)$ kg;
Question \ref{sec:c7-4} uses $(15,21,27)$ min. Their statistics and errors in both tables are
based on these midpoint estimates because individual observations are unknown.
\end{minipage}

\begin{minipage}{\linewidth}
\textbf{Errors relative to theory}

\par\smallskip
\begin{tabularx}{\linewidth}{@{}l >{\raggedright\arraybackslash}X r r r r@{}}
\toprule
Question & Statistic & Sample value & Theoretical value & $E$ & $E\%$ \\
\midrule
\ref{sec:c7-1} & $\bar{x}$ vs.\ $\mu$ (kg) & 8.167 & 9.000 & 0.833 & 9.26\% \\
 & $s^2$ vs.\ $\sigma^2$ (kg$^2$) & 5.500 & 7.000 & 1.500 & 21.43\% \\
 & $s$ vs.\ $\sigma$ (kg) & 2.345 & $\sqrt{7}\approx2.646$ & 0.301 & 11.36\% \\
\midrule
\ref{sec:c7-2} & $\bar{x}$ vs.\ $\mu$ ($^\circ$C) & 18.000 & 19.000 & 1.000 & 5.26\% \\
 & $s^2$ vs.\ $\sigma^2$ ($(^{\circ}\mathrm{C})^2$) & 8.500 & 6.000 & 2.500 & 41.67\% \\
 & $s$ vs.\ $\sigma$ ($^\circ$C) & 2.915 & $\sqrt{6}\approx2.449$ & 0.466 & 19.02\% \\
\midrule
\ref{sec:c7-3} & $\bar{x}$ vs.\ $\mu$ (parcels) & 7.444 & 7.000 & 0.444 & 6.35\% \\
 & $s^2$ vs.\ $\sigma^2$ (parcels$^2$) & 9.778 & 9.000 & 0.778 & 8.64\% \\
 & $s$ vs.\ $\sigma$ (parcels) & 3.127 & 3.000 & 0.127 & 4.23\% \\
\midrule
\ref{sec:c7-4} & $\bar{x}$ vs.\ $\mu$ (min) & 22.333 & 23.000 & 0.667 & 2.90\% \\
 & $s^2$ vs.\ $\sigma^2$ (min$^2$) & 25.000 & 30.000 & 5.000 & 16.67\% \\
 & $s$ vs.\ $\sigma$ (min) & 5.000 & $\sqrt{30}\approx5.477$ & 0.477 & 8.71\% \\
\bottomrule
\end{tabularx}
\par\smallskip

Sample values, theoretical values and absolute errors $E$ have the units shown
in the Statistic column; percentage errors are unitless.
\end{minipage}
\par
\endgroup
}

\end{document}
